Die Umlaufzeit ist umgekehrt proportional zur Winkelgeschwindigkeit, $T = 2\pi/\omega$. Mit $\omega' = (1+p)\,\omega$ und $p = \pzO$ wird
\begin{align}
 \frac{T'}{T} &= \frac{\omega}{\omega'} = \fTF = \fT
\end{align}
Die relative Änderung der Umlaufzeit ist
\begin{align}
 \frac{T'-T}{T} &= \dTF \\
   &= \dT \approx \result{\dTP{0}{2}}
\end{align}
Die Umlaufzeit wird also um knapp \pzO\ kürzer.
